\documentclass[11pt,a4paper]{article}
\usepackage[margin=1.7cm,top=1.5cm,bottom=1.5cm]{geometry}
\usepackage[fontset=fandol]{ctex}
\usepackage{amsmath,amssymb,graphicx,pgfplots,caption,booktabs}
\pgfplotsset{compat=1.18}\usepgfplotslibrary{groupplots}
\pgfplotsset{table/search path={figs/}}
\captionsetup{font=small,labelfont=bf}
\renewcommand{\thefigure}{S\arabic{figure}}
\setlength{\parindent}{2em}
\begin{document}
\noindent{\large\bfseries 八股条件背景模型：2019年测试段：每日背景更新}\par\medskip
沿用8:1:1划分后的历史测试段，各股68--69个保留交易日。蓝线只使用此前保留日的收益更新；参数及5日半衰期均已在训练、验证后固定。
\begin{center}\resizebox{.98\linewidth}{!}{\begin{tikzpicture}
\begin{groupplot}[group style={group size=2 by 4,horizontal sep=1.05cm,vertical sep=1.45cm},
width=7.45cm,height=4.45cm,
grid=major,grid style={gray!15},unbounded coords=jump,
tick label style={font=\scriptsize},title style={font=\small},label style={font=\small},
legend style={font=\footnotesize,draw=none,at={(1.12,1.23)},anchor=south,legend columns=4,column sep=.6em},
legend cell align=left,
ymin=0,ylabel={波动强度},scaled y ticks=false]
\nextgroupplot[title={(a) 中国平安 601318.SH},xmin=0,xmax=68,ymax=1.0800000000000001,xtick={0,17,34,51,68},xticklabels={09-11,10-14,11-06,11-29,12-24}]
\addplot[gray!70!black,line width=.5pt,no marks] table[x=day_index,y=R_d,col sep=comma]{eight_conditional_601318_test_scale.csv};
\addlegendentry{收盘后观测 $R_d$}
\addplot[blue!70!black,line width=.8pt,no marks] table[x=day_index,y=predicted_variance,col sep=comma]{eight_conditional_601318_test_scale.csv};
\addlegendentry{开盘前预报 $cq_d$}
\addplot[orange!85!black,dashed,line width=.8pt] coordinates {(0,1) (68,1)};
\addlegendentry{固定背景 $v_0=1$}
\nextgroupplot[title={(b) 贵州茅台 600519.SH},xmin=0,xmax=68,ymax=1.8954455083495421,xtick={0,17,34,51,68},xticklabels={09-11,10-14,11-06,11-29,12-24},ylabel={}]
\addplot[gray!70!black,line width=.5pt,no marks] table[x=day_index,y=R_d,col sep=comma]{eight_conditional_600519_test_scale.csv};
\addplot[blue!70!black,line width=.8pt,no marks] table[x=day_index,y=predicted_variance,col sep=comma]{eight_conditional_600519_test_scale.csv};
\addplot[orange!85!black,dashed,line width=.8pt] coordinates {(0,1) (68,1)};
\nextgroupplot[title={(c) 中信证券 600030.SH},xmin=0,xmax=67,ymax=1.8342792018575214,xtick={0,16,33,50,67},xticklabels={09-11,10-11,11-06,11-29,12-24}]
\addplot[gray!70!black,line width=.5pt,no marks] table[x=day_index,y=R_d,col sep=comma]{eight_conditional_600030_test_scale.csv};
\addplot[blue!70!black,line width=.8pt,no marks] table[x=day_index,y=predicted_variance,col sep=comma]{eight_conditional_600030_test_scale.csv};
\addplot[orange!85!black,dashed,line width=.8pt] coordinates {(0,1) (67,1)};
\nextgroupplot[title={(d) 伊利股份 600887.SH},xmin=0,xmax=68,ymax=1.5364969817666234,xtick={0,17,34,51,68},xticklabels={09-11,10-14,11-06,11-29,12-24},ylabel={}]
\addplot[gray!70!black,line width=.5pt,no marks] table[x=day_index,y=R_d,col sep=comma]{eight_conditional_600887_test_scale.csv};
\addplot[blue!70!black,line width=.8pt,no marks] table[x=day_index,y=predicted_variance,col sep=comma]{eight_conditional_600887_test_scale.csv};
\addplot[orange!85!black,dashed,line width=.8pt] coordinates {(0,1) (68,1)};
\nextgroupplot[title={(e) 招商银行 600036.SH},xmin=0,xmax=68,ymax=1.3295361690930771,xtick={0,17,34,51,68},xticklabels={09-11,10-14,11-06,11-29,12-24}]
\addplot[gray!70!black,line width=.5pt,no marks] table[x=day_index,y=R_d,col sep=comma]{eight_conditional_600036_test_scale.csv};
\addplot[blue!70!black,line width=.8pt,no marks] table[x=day_index,y=predicted_variance,col sep=comma]{eight_conditional_600036_test_scale.csv};
\addplot[orange!85!black,dashed,line width=.8pt] coordinates {(0,1) (68,1)};
\nextgroupplot[title={(f) 华友钴业 603799.SH},xmin=0,xmax=68,ymax=2.8005366175905007,xtick={0,17,34,51,68},xticklabels={09-11,10-14,11-06,11-29,12-24},ylabel={}]
\addplot[gray!70!black,line width=.5pt,no marks] table[x=day_index,y=R_d,col sep=comma]{eight_conditional_603799_test_scale.csv};
\addplot[blue!70!black,line width=.8pt,no marks] table[x=day_index,y=predicted_variance,col sep=comma]{eight_conditional_603799_test_scale.csv};
\addplot[orange!85!black,dashed,line width=.8pt] coordinates {(0,1) (68,1)};
\nextgroupplot[title={(g) 方大炭素 600516.SH},xmin=0,xmax=68,ymax=3.0746623306574792,xtick={0,17,34,51,68},xticklabels={09-11,10-14,11-06,11-29,12-24},xlabel={观测交易日}]
\addplot[gray!70!black,line width=.5pt,no marks] table[x=day_index,y=R_d,col sep=comma]{eight_conditional_600516_test_scale.csv};
\addplot[blue!70!black,line width=.8pt,no marks] table[x=day_index,y=predicted_variance,col sep=comma]{eight_conditional_600516_test_scale.csv};
\addplot[orange!85!black,dashed,line width=.8pt] coordinates {(0,1) (68,1)};
\nextgroupplot[title={(h) 中国建筑 601668.SH},xmin=0,xmax=68,ymax=1.7688298790030172,xtick={0,17,34,51,68},xticklabels={09-11,10-14,11-06,11-29,12-24},ylabel={},xlabel={观测交易日}]
\addplot[gray!70!black,line width=.5pt,no marks] table[x=day_index,y=R_d,col sep=comma]{eight_conditional_601668_test_scale.csv};
\addplot[blue!70!black,line width=.8pt,no marks] table[x=day_index,y=predicted_variance,col sep=comma]{eight_conditional_601668_test_scale.csv};
\addplot[orange!85!black,dashed,line width=.8pt] coordinates {(0,1) (68,1)};
\end{groupplot}\end{tikzpicture}
}
\captionof{figure}{灰线为收盘后计算的$R_d$，蓝线为当天开盘前已知的$cq_d$，橙虚线为$v_0=1$。日序号按实际观测日排列，标签为月--日；休市间隔压缩。纵轴以各股训练方差为单位，面板范围可不同。}\end{center}

\clearpage
\noindent{\large\bfseries 八股条件背景模型：2026年：每日背景更新}\par\medskip
各股从$q_1=1$重新启动；华友钴业覆盖6月30日至8月31日，共45日，其余7股均为160日。短样本的起始状态影响更大，不能将其与完整年份等量看待。
\begin{center}\resizebox{.98\linewidth}{!}{\begin{tikzpicture}
\begin{groupplot}[group style={group size=2 by 4,horizontal sep=1.05cm,vertical sep=1.45cm},
width=7.45cm,height=4.45cm,
grid=major,grid style={gray!15},unbounded coords=jump,
tick label style={font=\scriptsize},title style={font=\small},label style={font=\small},
legend style={font=\footnotesize,draw=none,at={(1.12,1.23)},anchor=south,legend columns=4,column sep=.6em},
legend cell align=left,
ymin=0,ylabel={波动强度},scaled y ticks=false]
\nextgroupplot[title={(a) 中国平安 601318.SH},xmin=0,xmax=159,ymax=3.9257840446220253,xtick={0,39,79,119,159},xticklabels={01-05,03-09,05-08,07-06,08-31}]
\addplot[gray!70!black,line width=.5pt,no marks] table[x=day_index,y=R_d,col sep=comma]{eight_conditional_601318_2026_scale.csv};
\addlegendentry{收盘后观测 $R_d$}
\addplot[blue!70!black,line width=.8pt,no marks] table[x=day_index,y=predicted_variance,col sep=comma]{eight_conditional_601318_2026_scale.csv};
\addlegendentry{开盘前预报 $cq_d$}
\addplot[orange!85!black,dashed,line width=.8pt] coordinates {(0,1) (159,1)};
\addlegendentry{固定背景 $v_0=1$}
\nextgroupplot[title={(b) 贵州茅台 600519.SH},xmin=0,xmax=159,ymax=4.244626178714654,xtick={0,39,79,119,159},xticklabels={01-05,03-09,05-08,07-06,08-31},ylabel={}]
\addplot[gray!70!black,line width=.5pt,no marks] table[x=day_index,y=R_d,col sep=comma]{eight_conditional_600519_2026_scale.csv};
\addplot[blue!70!black,line width=.8pt,no marks] table[x=day_index,y=predicted_variance,col sep=comma]{eight_conditional_600519_2026_scale.csv};
\addplot[orange!85!black,dashed,line width=.8pt] coordinates {(0,1) (159,1)};
\nextgroupplot[title={(c) 中信证券 600030.SH},xmin=0,xmax=159,ymax=3.9703961348903887,xtick={0,39,79,119,159},xticklabels={01-05,03-09,05-08,07-06,08-31}]
\addplot[gray!70!black,line width=.5pt,no marks] table[x=day_index,y=R_d,col sep=comma]{eight_conditional_600030_2026_scale.csv};
\addplot[blue!70!black,line width=.8pt,no marks] table[x=day_index,y=predicted_variance,col sep=comma]{eight_conditional_600030_2026_scale.csv};
\addplot[orange!85!black,dashed,line width=.8pt] coordinates {(0,1) (159,1)};
\nextgroupplot[title={(d) 伊利股份 600887.SH},xmin=0,xmax=159,ymax=1.2962537431017505,xtick={0,39,79,119,159},xticklabels={01-05,03-09,05-08,07-06,08-31},ylabel={}]
\addplot[gray!70!black,line width=.5pt,no marks] table[x=day_index,y=R_d,col sep=comma]{eight_conditional_600887_2026_scale.csv};
\addplot[blue!70!black,line width=.8pt,no marks] table[x=day_index,y=predicted_variance,col sep=comma]{eight_conditional_600887_2026_scale.csv};
\addplot[orange!85!black,dashed,line width=.8pt] coordinates {(0,1) (159,1)};
\nextgroupplot[title={(e) 招商银行 600036.SH},xmin=0,xmax=159,ymax=1.3153512499406812,xtick={0,39,79,119,159},xticklabels={01-05,03-09,05-08,07-06,08-31}]
\addplot[gray!70!black,line width=.5pt,no marks] table[x=day_index,y=R_d,col sep=comma]{eight_conditional_600036_2026_scale.csv};
\addplot[blue!70!black,line width=.8pt,no marks] table[x=day_index,y=predicted_variance,col sep=comma]{eight_conditional_600036_2026_scale.csv};
\addplot[orange!85!black,dashed,line width=.8pt] coordinates {(0,1) (159,1)};
\nextgroupplot[title={(f) 华友钴业 603799.SH（45日）},xmin=0,xmax=44,ymax=2.7969500711688426,xtick={0,11,22,33,44},xticklabels={06-30,07-15,07-30,08-14,08-31},ylabel={}]
\addplot[gray!70!black,line width=.5pt,no marks] table[x=day_index,y=R_d,col sep=comma]{eight_conditional_603799_2026_scale.csv};
\addplot[blue!70!black,line width=.8pt,no marks] table[x=day_index,y=predicted_variance,col sep=comma]{eight_conditional_603799_2026_scale.csv};
\addplot[orange!85!black,dashed,line width=.8pt] coordinates {(0,1) (44,1)};
\nextgroupplot[title={(g) 方大炭素 600516.SH},xmin=0,xmax=159,ymax=9.5202236523666102,xtick={0,39,79,119,159},xticklabels={01-05,03-09,05-08,07-06,08-31},xlabel={观测交易日}]
\addplot[gray!70!black,line width=.5pt,no marks] table[x=day_index,y=R_d,col sep=comma]{eight_conditional_600516_2026_scale.csv};
\addplot[blue!70!black,line width=.8pt,no marks] table[x=day_index,y=predicted_variance,col sep=comma]{eight_conditional_600516_2026_scale.csv};
\addplot[orange!85!black,dashed,line width=.8pt] coordinates {(0,1) (159,1)};
\nextgroupplot[title={(h) 中国建筑 601668.SH},xmin=0,xmax=159,ymax=1.7473073089947073,xtick={0,39,79,119,159},xticklabels={01-05,03-09,05-08,07-06,08-31},ylabel={},xlabel={观测交易日}]
\addplot[gray!70!black,line width=.5pt,no marks] table[x=day_index,y=R_d,col sep=comma]{eight_conditional_601668_2026_scale.csv};
\addplot[blue!70!black,line width=.8pt,no marks] table[x=day_index,y=predicted_variance,col sep=comma]{eight_conditional_601668_2026_scale.csv};
\addplot[orange!85!black,dashed,line width=.8pt] coordinates {(0,1) (159,1)};
\end{groupplot}\end{tikzpicture}
}
\captionof{figure}{线条含义同图S1。$R_d$是当日波动强度的观测估计，并非真实隐含方差；蓝线在突发波动时可能滞后。所有股票共用5日半衰期，形状参数和尺度校准系数各自由训练段估计。}\end{center}

\clearpage
\noindent{\large\bfseries 八股条件背景模型：2019年测试段：经验密度与预测密度}\par\medskip
每股沿用其训练段的80箱边界，经验密度以测试段全部有效收益数归一，空箱不绘制。预测曲线没有用测试期标准差重新缩放。
\begin{center}\resizebox{.98\linewidth}{!}{\begin{tikzpicture}
\begin{groupplot}[group style={group size=2 by 4,horizontal sep=1.05cm,vertical sep=1.45cm},
width=7.45cm,height=4.45cm,
grid=major,grid style={gray!15},unbounded coords=jump,
tick label style={font=\scriptsize},title style={font=\small},label style={font=\small},
legend style={font=\footnotesize,draw=none,at={(1.12,1.23)},anchor=south,legend columns=4,column sep=.6em},
legend cell align=left,
ymode=log,ymin=0.00001,ymax=5,ylabel={概率密度},ytick={0.00001,0.001,0.1,1}]
\nextgroupplot[title={(a) 中国平安 601318.SH},xmin=-4.595921659771838,xmax=4.595921659771838]
\addplot[black!70,only marks,mark=*,mark size=.85pt] table[x=x,y=plot_density,col sep=comma]{eight_conditional_601318_test_hist.csv};
\addlegendentry{经验密度}
\addplot[orange!85!black,line width=.8pt] table[x=x,y=static_N8,col sep=comma]{eight_conditional_601318_test_density.csv};
\addlegendentry{静态似然八层}
\addplot[blue!70!black,line width=.8pt] table[x=x,y=dynamic_N8,col sep=comma]{eight_conditional_601318_test_density.csv};
\addlegendentry{条件八层}
\addplot[violet!85!black,dashdotted,line width=.8pt] table[x=x,y=dynamic_Gaussian,col sep=comma]{eight_conditional_601318_test_density.csv};
\addlegendentry{同规则高斯}
\nextgroupplot[title={(b) 贵州茅台 600519.SH},xmin=-5.0124536266210145,xmax=5.0124536266210145,ylabel={}]
\addplot[black!70,only marks,mark=*,mark size=.85pt] table[x=x,y=plot_density,col sep=comma]{eight_conditional_600519_test_hist.csv};
\addplot[orange!85!black,line width=.8pt] table[x=x,y=static_N8,col sep=comma]{eight_conditional_600519_test_density.csv};
\addplot[blue!70!black,line width=.8pt] table[x=x,y=dynamic_N8,col sep=comma]{eight_conditional_600519_test_density.csv};
\addplot[violet!85!black,dashdotted,line width=.8pt] table[x=x,y=dynamic_Gaussian,col sep=comma]{eight_conditional_600519_test_density.csv};
\nextgroupplot[title={(c) 中信证券 600030.SH},xmin=-5.660159476299526,xmax=5.660159476299526]
\addplot[black!70,only marks,mark=*,mark size=.85pt] table[x=x,y=plot_density,col sep=comma]{eight_conditional_600030_test_hist.csv};
\addplot[orange!85!black,line width=.8pt] table[x=x,y=static_N8,col sep=comma]{eight_conditional_600030_test_density.csv};
\addplot[blue!70!black,line width=.8pt] table[x=x,y=dynamic_N8,col sep=comma]{eight_conditional_600030_test_density.csv};
\addplot[violet!85!black,dashdotted,line width=.8pt] table[x=x,y=dynamic_Gaussian,col sep=comma]{eight_conditional_600030_test_density.csv};
\nextgroupplot[title={(d) 伊利股份 600887.SH},xmin=-4.8198806302473765,xmax=4.8198806302473765,ylabel={}]
\addplot[black!70,only marks,mark=*,mark size=.85pt] table[x=x,y=plot_density,col sep=comma]{eight_conditional_600887_test_hist.csv};
\addplot[orange!85!black,line width=.8pt] table[x=x,y=static_N8,col sep=comma]{eight_conditional_600887_test_density.csv};
\addplot[blue!70!black,line width=.8pt] table[x=x,y=dynamic_N8,col sep=comma]{eight_conditional_600887_test_density.csv};
\addplot[violet!85!black,dashdotted,line width=.8pt] table[x=x,y=dynamic_Gaussian,col sep=comma]{eight_conditional_600887_test_density.csv};
\nextgroupplot[title={(e) 招商银行 600036.SH},xmin=-4.728278701608616,xmax=4.728278701608616]
\addplot[black!70,only marks,mark=*,mark size=.85pt] table[x=x,y=plot_density,col sep=comma]{eight_conditional_600036_test_hist.csv};
\addplot[orange!85!black,line width=.8pt] table[x=x,y=static_N8,col sep=comma]{eight_conditional_600036_test_density.csv};
\addplot[blue!70!black,line width=.8pt] table[x=x,y=dynamic_N8,col sep=comma]{eight_conditional_600036_test_density.csv};
\addplot[violet!85!black,dashdotted,line width=.8pt] table[x=x,y=dynamic_Gaussian,col sep=comma]{eight_conditional_600036_test_density.csv};
\nextgroupplot[title={(f) 华友钴业 603799.SH},xmin=-5.120680785735817,xmax=5.120680785735817,ylabel={}]
\addplot[black!70,only marks,mark=*,mark size=.85pt] table[x=x,y=plot_density,col sep=comma]{eight_conditional_603799_test_hist.csv};
\addplot[orange!85!black,line width=.8pt] table[x=x,y=static_N8,col sep=comma]{eight_conditional_603799_test_density.csv};
\addplot[blue!70!black,line width=.8pt] table[x=x,y=dynamic_N8,col sep=comma]{eight_conditional_603799_test_density.csv};
\addplot[violet!85!black,dashdotted,line width=.8pt] table[x=x,y=dynamic_Gaussian,col sep=comma]{eight_conditional_603799_test_density.csv};
\nextgroupplot[title={(g) 方大炭素 600516.SH},xmin=-6.00372101828073,xmax=6.00372101828073,xlabel={归一化收益 $x$}]
\addplot[black!70,only marks,mark=*,mark size=.85pt] table[x=x,y=plot_density,col sep=comma]{eight_conditional_600516_test_hist.csv};
\addplot[orange!85!black,line width=.8pt] table[x=x,y=static_N8,col sep=comma]{eight_conditional_600516_test_density.csv};
\addplot[blue!70!black,line width=.8pt] table[x=x,y=dynamic_N8,col sep=comma]{eight_conditional_600516_test_density.csv};
\addplot[violet!85!black,dashdotted,line width=.8pt] table[x=x,y=dynamic_Gaussian,col sep=comma]{eight_conditional_600516_test_density.csv};
\nextgroupplot[title={(h) 中国建筑 601668.SH},xmin=-5.1961255383408735,xmax=5.1961255383408735,ylabel={},xlabel={归一化收益 $x$}]
\addplot[black!70,only marks,mark=*,mark size=.85pt] table[x=x,y=plot_density,col sep=comma]{eight_conditional_601668_test_hist.csv};
\addplot[orange!85!black,line width=.8pt] table[x=x,y=static_N8,col sep=comma]{eight_conditional_601668_test_density.csv};
\addplot[blue!70!black,line width=.8pt] table[x=x,y=dynamic_N8,col sep=comma]{eight_conditional_601668_test_density.csv};
\addplot[violet!85!black,dashdotted,line width=.8pt] table[x=x,y=dynamic_Gaussian,col sep=comma]{eight_conditional_601668_test_density.csv};
\end{groupplot}\end{tikzpicture}
}
\captionof{figure}{黑点为经验密度，橙线为静态似然八层，蓝线为条件八层，紫线为同规则高斯。条件曲线按每日有效收益数加权，仅用于展示；逐点评分仍使用收益发生前确定的条件密度。各面板的横轴范围沿用各股训练分箱。}\end{center}

\clearpage
\noindent{\large\bfseries 八股条件背景模型：2026年：经验密度与预测密度}\par\medskip
使用2017--2019年训练、验证后固定的参数。华友钴业只有45日、2070条收益，其余7股均有160日、7360条收益；尾部散点的疏密受到样本量影响。
\begin{center}\resizebox{.98\linewidth}{!}{\begin{tikzpicture}
\begin{groupplot}[group style={group size=2 by 4,horizontal sep=1.05cm,vertical sep=1.45cm},
width=7.45cm,height=4.45cm,
grid=major,grid style={gray!15},unbounded coords=jump,
tick label style={font=\scriptsize},title style={font=\small},label style={font=\small},
legend style={font=\footnotesize,draw=none,at={(1.12,1.23)},anchor=south,legend columns=4,column sep=.6em},
legend cell align=left,
ymode=log,ymin=0.00001,ymax=5,ylabel={概率密度},ytick={0.00001,0.001,0.1,1}]
\nextgroupplot[title={(a) 中国平安 601318.SH},xmin=-4.595921659771838,xmax=4.595921659771838]
\addplot[black!70,only marks,mark=*,mark size=.85pt] table[x=x,y=plot_density,col sep=comma]{eight_conditional_601318_2026_hist.csv};
\addlegendentry{经验密度}
\addplot[orange!85!black,line width=.8pt] table[x=x,y=static_N8,col sep=comma]{eight_conditional_601318_2026_density.csv};
\addlegendentry{静态似然八层}
\addplot[blue!70!black,line width=.8pt] table[x=x,y=dynamic_N8,col sep=comma]{eight_conditional_601318_2026_density.csv};
\addlegendentry{条件八层}
\addplot[violet!85!black,dashdotted,line width=.8pt] table[x=x,y=dynamic_Gaussian,col sep=comma]{eight_conditional_601318_2026_density.csv};
\addlegendentry{同规则高斯}
\nextgroupplot[title={(b) 贵州茅台 600519.SH},xmin=-5.0124536266210145,xmax=5.0124536266210145,ylabel={}]
\addplot[black!70,only marks,mark=*,mark size=.85pt] table[x=x,y=plot_density,col sep=comma]{eight_conditional_600519_2026_hist.csv};
\addplot[orange!85!black,line width=.8pt] table[x=x,y=static_N8,col sep=comma]{eight_conditional_600519_2026_density.csv};
\addplot[blue!70!black,line width=.8pt] table[x=x,y=dynamic_N8,col sep=comma]{eight_conditional_600519_2026_density.csv};
\addplot[violet!85!black,dashdotted,line width=.8pt] table[x=x,y=dynamic_Gaussian,col sep=comma]{eight_conditional_600519_2026_density.csv};
\nextgroupplot[title={(c) 中信证券 600030.SH},xmin=-5.660159476299526,xmax=5.660159476299526]
\addplot[black!70,only marks,mark=*,mark size=.85pt] table[x=x,y=plot_density,col sep=comma]{eight_conditional_600030_2026_hist.csv};
\addplot[orange!85!black,line width=.8pt] table[x=x,y=static_N8,col sep=comma]{eight_conditional_600030_2026_density.csv};
\addplot[blue!70!black,line width=.8pt] table[x=x,y=dynamic_N8,col sep=comma]{eight_conditional_600030_2026_density.csv};
\addplot[violet!85!black,dashdotted,line width=.8pt] table[x=x,y=dynamic_Gaussian,col sep=comma]{eight_conditional_600030_2026_density.csv};
\nextgroupplot[title={(d) 伊利股份 600887.SH},xmin=-4.8198806302473765,xmax=4.8198806302473765,ylabel={}]
\addplot[black!70,only marks,mark=*,mark size=.85pt] table[x=x,y=plot_density,col sep=comma]{eight_conditional_600887_2026_hist.csv};
\addplot[orange!85!black,line width=.8pt] table[x=x,y=static_N8,col sep=comma]{eight_conditional_600887_2026_density.csv};
\addplot[blue!70!black,line width=.8pt] table[x=x,y=dynamic_N8,col sep=comma]{eight_conditional_600887_2026_density.csv};
\addplot[violet!85!black,dashdotted,line width=.8pt] table[x=x,y=dynamic_Gaussian,col sep=comma]{eight_conditional_600887_2026_density.csv};
\nextgroupplot[title={(e) 招商银行 600036.SH},xmin=-4.728278701608616,xmax=4.728278701608616]
\addplot[black!70,only marks,mark=*,mark size=.85pt] table[x=x,y=plot_density,col sep=comma]{eight_conditional_600036_2026_hist.csv};
\addplot[orange!85!black,line width=.8pt] table[x=x,y=static_N8,col sep=comma]{eight_conditional_600036_2026_density.csv};
\addplot[blue!70!black,line width=.8pt] table[x=x,y=dynamic_N8,col sep=comma]{eight_conditional_600036_2026_density.csv};
\addplot[violet!85!black,dashdotted,line width=.8pt] table[x=x,y=dynamic_Gaussian,col sep=comma]{eight_conditional_600036_2026_density.csv};
\nextgroupplot[title={(f) 华友钴业 603799.SH（45日）},xmin=-5.120680785735817,xmax=5.120680785735817,ylabel={}]
\addplot[black!70,only marks,mark=*,mark size=.85pt] table[x=x,y=plot_density,col sep=comma]{eight_conditional_603799_2026_hist.csv};
\addplot[orange!85!black,line width=.8pt] table[x=x,y=static_N8,col sep=comma]{eight_conditional_603799_2026_density.csv};
\addplot[blue!70!black,line width=.8pt] table[x=x,y=dynamic_N8,col sep=comma]{eight_conditional_603799_2026_density.csv};
\addplot[violet!85!black,dashdotted,line width=.8pt] table[x=x,y=dynamic_Gaussian,col sep=comma]{eight_conditional_603799_2026_density.csv};
\nextgroupplot[title={(g) 方大炭素 600516.SH},xmin=-6.00372101828073,xmax=6.00372101828073,xlabel={归一化收益 $x$}]
\addplot[black!70,only marks,mark=*,mark size=.85pt] table[x=x,y=plot_density,col sep=comma]{eight_conditional_600516_2026_hist.csv};
\addplot[orange!85!black,line width=.8pt] table[x=x,y=static_N8,col sep=comma]{eight_conditional_600516_2026_density.csv};
\addplot[blue!70!black,line width=.8pt] table[x=x,y=dynamic_N8,col sep=comma]{eight_conditional_600516_2026_density.csv};
\addplot[violet!85!black,dashdotted,line width=.8pt] table[x=x,y=dynamic_Gaussian,col sep=comma]{eight_conditional_600516_2026_density.csv};
\nextgroupplot[title={(h) 中国建筑 601668.SH},xmin=-5.1961255383408735,xmax=5.1961255383408735,ylabel={},xlabel={归一化收益 $x$}]
\addplot[black!70,only marks,mark=*,mark size=.85pt] table[x=x,y=plot_density,col sep=comma]{eight_conditional_601668_2026_hist.csv};
\addplot[orange!85!black,line width=.8pt] table[x=x,y=static_N8,col sep=comma]{eight_conditional_601668_2026_density.csv};
\addplot[blue!70!black,line width=.8pt] table[x=x,y=dynamic_N8,col sep=comma]{eight_conditional_601668_2026_density.csv};
\addplot[violet!85!black,dashdotted,line width=.8pt] table[x=x,y=dynamic_Gaussian,col sep=comma]{eight_conditional_601668_2026_density.csv};
\end{groupplot}\end{tikzpicture}
}
\captionof{figure}{线条与归一化规则同图S3。静态与条件八层使用相同训练样本；高斯对照共享背景更新规则，其尺度校准系数单独从训练段估计。零收益集中及尾部偏差仍可能存在，曲线接近不能替代评分与覆盖率检验。}\end{center}

\end{document}
